% ============================================================
\section{Funtor $P: Fin \to Sub$}
\label{sec:functor-fin-sub}

\subsection{Descrição das Categorias}

\begin{definition}[Categoria $\mathbf{Sub}$ (modelo MATLAB/Octave)]
Um objeto é um vetor de inteiros (os elementos do conjunto). Um morfismo $A \to B$ existe apenas se $A \subseteq B$, e é representado pelo vetor lógico de inclusão.
\end{definition}

\begin{lstlisting}[style=matlab, caption={Morfismo de inclusão em $\mathbf{Sub}$}]
A = [1, 3];  B = [1, 2, 3, 4];
is_sub = all(ismember(A, B));  % verifica se A esta contido em B
incl   = ismember(B, A);       % vector logico da inclusao
\end{lstlisting}

\subsection{Funtor Covariante $P: Fin \to Sub$ (Conjunto das Partes)}

O \emph{funtor potência} $P: \mathbf{Fin} \to \mathbf{Sub}$ envia cada conjunto finito $n$ ao conjunto das suas partes $\mathcal{P}(n) = \{S \subseteq \{1,\ldots,n\}\}$, e cada aplicação total $f: n \to m$ à aplicação
\[
    P(f): \mathcal{P}(n) \to \mathcal{P}(m), \quad S \;\longmapsto\; f(S) = \{f(i) : i \in S\}.
\]

\begin{lstlisting}[style=matlab, caption={Funtor potência $P$ e imagem direta}]
function PS = power_set(n)
    PS = {};
    for k = 0:2^n-1
        PS{end+1} = find(bitand(k, 2.^(0:n-1)));
    end
end

function img = direct_image(f, S)
    img = unique(f(S));  % imagem directa de S por f
end
\end{lstlisting}

\begin{remark}
Este funtor preserva inclusões: se $S \subseteq T$ então $f(S) \subseteq f(T)$, pelo que está bem definido em $\mathbf{Sub}$ como codomínio.
\end{remark}

\subsection{Funtor de Esquecimento $U: Sub \to Fin$}

O funtor de esquecimento $U: \mathbf{Sub} \to \mathbf{Fin}$ envia cada objeto de $\mathbf{Sub}$ (um subconjunto $A \subseteq B$) ao conjunto finito $A$, ignorando a estrutura de inclusão, e cada morfismo (inclusão $A \hookrightarrow B$) à correspondente aplicação total em $\mathbf{Fin}$.

\subsection{Funtor Contravariante $P^*: Fin^{\mathrm{op}} \to Sub$ (Imagem Inversa)}

O \emph{funtor imagem inversa} $P^*: \mathbf{Fin}^{\mathrm{op}} \to \mathbf{Sub}$ envia cada aplicação $f: n \to m$ à aplicação
\[
    P^*(f): \mathcal{P}(m) \to \mathcal{P}(n), \quad T \;\longmapsto\; f^{-1}(T) = \{i \in n : f(i) \in T\}.
\]
Este funtor é \emph{contravariante} porque inverte a direção dos morfismos.

\begin{lstlisting}[style=matlab, caption={Funtor contravariante $P^*$ (pré-imagem)}]
function preimg = inverse_image(f, T)
    preimg = find(ismember(f, T));  % pre-imagem de T por f
end
\end{lstlisting}

\subsection{Transformações Naturais entre Funtores $Fin \to Sub$}

Dados dois funtores $F, G: \mathbf{Fin} \to \mathbf{Sub}$, uma transformação natural $\eta: F \Rightarrow G$ consiste numa família de morfismos em $\mathbf{Sub}$,
\[
    \eta_n: F(n) \hookrightarrow G(n) \quad \text{(i.e., } F(n) \subseteq G(n) \text{ para todo } n\text{)},
\]
tal que para todo morfismo $f: n \to m$ em $\mathbf{Fin}$:
\[
    G(f)\bigl(\eta_n(x)\bigr) \;=\; \eta_m\bigl(F(f)(x)\bigr), \quad \forall\, x \in F(n).
\]

\begin{remark}
Seja $F = P$ (funtor potência) e $G$ o funtor que associa a $n$ o conjunto de todos os multiconjuntos de $\{1,\ldots,n\}$. A inclusão natural de subconjuntos em multiconjuntos define uma transformação natural $\eta: P \Rightarrow G$.
\end{remark}
